\chapter{Introduction}

Large language model (LLM) evaluations are usually easiest to run through APIs.
Programmatic access makes it possible to fix prompts, set sampling parameters,
log responses, and repeat collections in a controlled way. This is the logic
behind benchmark suites such as HELM by Liang et al.\ \cite*{helm}, which
emphasize standardized scenarios, metrics, and collection conditions. But most end users do not
encounter LLMs through bare evaluation endpoints. They encounter productized
applications with interface defaults, source displays, tool use, personalization,
moderation messages, and session context. For audit work, those layers are not
just presentation: they can change the answer a user receives.

This thesis treats the deployment surface as part of the system under audit.
That choice follows the broader algorithm-auditing tradition described by
Sandvig et al.\ \cite*{sandvig_auditing}, where controlled inputs are used to
study opaque deployed systems from their observable outputs. It also follows
Raji et al.'s accountability work \cite*{raji_gap}, which argues that audits
must document the system boundary and the conditions under which evidence is
collected. This framing is especially important for foundation-model systems
because their risks are sociotechnical and deployment-dependent, not only
properties of a model checkpoint, as discussed by Bommasani et al.\
\cite*{foundation_models} and Weidinger et al.\ \cite*{weidinger_risks}.

The closest related access-surface study is Dual Standards by Lipphardt et al.\
\cite*{dualstandards}, which compares moderation behavior between API and WebUI
interfaces for LLMs. Its relevance is not that the same direction of effect
should appear here. Rather, it shows that API and WebUI surfaces can be
empirically non-equivalent. This thesis asks the same kind of access-surface
question in a narrower setting: a bare OpenAI API condition compared with the
official ChatGPT app on one frozen diagnostic prompt bank.

The research question is:

\begin{quote}
How do observable outputs differ between a bare OpenAI API condition and the
official productized ChatGPT app deployment on a frozen diagnostic prompt suite?
\end{quote}

The study uses a 60-prompt bank covering allowed-sensitive over-refusal,
refusal-expected guardrailing, boundary safety framing, instruction following,
socially sensitive recommendations, and neutral factual controls. The core
comparison is the OpenAI Responses API using \texttt{gpt-4o} with no custom
system prompt versus screenshot-backed collection from the official ChatGPT
consumer app. Claude and Gemini API runs are included only as optional baselines.
The design deliberately keeps the claim narrow. It does not isolate the effect
of a hidden system prompt, retrieval policy, tool policy, routing decision, or
moderation layer. The evidence supports only observable differences between a
bare API condition and a productized app deployment under documented collection
conditions.

The empirical result is a compact diagnostic audit, not a population estimate.
After transcript correction and manual review, 17 of 60 API/app pairs show
substantive observable divergences. Ten have response-core anchors, mainly
refusal-granularity and factual-control differences. Seven are product-surface
or context cases, such as visible sources or localized information. Targeted app
repetitions for the divergent cases support most of the original case inventory
while keeping app-side stability claims scoped to that packet. Inter-rater
reliability also sharpens the interpretation: refusal status is reliable, while
soft safety-framing strength is treated only as supporting evidence.

The contribution is therefore methodological as much as empirical. The thesis
operationalizes deployment surface as an audit variable, builds a
benchmark-grounded prompt suite, applies a manually reviewed divergence taxonomy,
and preserves screenshot-backed evidence with correction and reliability
materials. Its central conclusion is cautious: for this prompt bank and
collection window, the bare API condition is not a neutral proxy for the official
ChatGPT app. LLM audits should therefore name and document the access surface
they evaluate.

Section 2 reviews related work. Section 3 describes the method. Section 4
reports results. Sections 5 to 8 discuss implications, limitations, future work,
and conclusions. The appendix provides the redacted prompt bank, coding guide,
correction audit, reliability materials, and artifact inventory.
